\chapter{Abstract}

P=NP

\chapter{Einleitung}

\section{Einführung in das Thema}

\section{Zielstellung}

Ziel der Arbeit ist es, Anwendungen, die das Microservice-Paradigma verwenden, mit Werkzeugen und Methoden der modellgetriebenen Softwareentwicklung zu entwickeln.

Dazu wird ein Metamodell konzipiert, welches die grundlegenden Eigenschaften des Paradigmas abstrahiert. Ergänzt wird dies durch die Berücksichtigung von Konzepten des Domain-Driven Design zur Modellierung der Geschäftslogik.

Die mithilfe dieses Metamodells generierten Anwendungen sollen in der Lage sein, bereits existierende, sogenannte Microservice Chassis – wiederverwendbare Grundbausteine eines Microservices – zu verwenden. Daraufhin wird erforscht, welche Möglichkeiten es gibt, diese generierten Anwendungen automatisiert, in Abhängigkeit von Versionsaktualisierungen der Chassis, zu migrieren.

Außerdem wird untersucht, inwiefern die der Anwendung zugrunde liegenden IT-Infrastruktur in solch einem Metamodell berücksichtigt werden kann und welche Möglichkeiten sich daraus ergeben.